% Converted from README.md
\section{Testiranje pomoću objekata imitatora (engl. Mock testing)}\label{testiranje-pomoux107u-objekata-imitatora-engl-mock-testing}

Prilikom pisanja jediničnih testova fokusiramo se samo na jednu funkciju i ispitujemo njeno ponašanje u kontrolisanom okruženju. Sadržaj testa je uvek skup inicijalizacija okruženja, pokretanje funkcije koju testiramo i zatim poređenje dobijenog i očekivanog rezultata. Nekada je kod takav da je nemoguće napraviti kontrolisano okruženje za testiranje, npr. korišćenje sistemskih poziva, podataka iz baze podataka, mrežna komunikacija i sl. U tim situacijama pribegava se pisanju klasa koje imitiraju realne objekte u svrhu testiranja.

Na primer, funkcija može da obrađuje podatke iz datoteke i uzima naziv datoteke kao ulazni parametar. Bolje rešenje je prepraviti funkciju tako da radi sa ulaznim tokom koji treba da obradi. Funkcija potom neće raditi sve stvari kao ranije - otvarati datoteku i obrađivati podatke.

U jediničnom testu, objekti imitatori mogu da imitiraju ponašanje kompleksnog stvarnog objekta. Korisni su u situacijama kada stvarni objekat nije praktično ili je nemoguće uklopiti u jedinični test. Objekat imitator se obično koristi ukoliko stvarni objekat ima neku od narednih karakteristika:

\begin{itemize}
\tightlist
\item
  obezbeđuje nedeterministički rezultat (npr. trenutno vreme, trenutnu temperaturu, ...)
\item
  ima stanja koja je teško kreirati ili reprodukovati, (npr. greška u mrežnoj komunikaciji)
\item
  spor je (npr. baza podataka, koja bi pre svakog testa morala biti inicijalizovana)
\item
  ne postoji još uvek, ili može promeniti ponašanje u budućnosti
\item
  morao bi da dobije nove informacije i metode da bismo mogli da ga koristimo za testiranje, a inače mu nisu potrebne
\end{itemize}

Objekti imitatori treba da imaju isti interfejs kao stvarni objekti koje imitiraju. Tako omogućavaju da objekat koji ih koristi ne pravi razliku između stvarnog objekta i imitatora. Mnogi radni okviri za objekte imitatore omogućavaju da se navede objekat koje klase se imitira, a zatim da programer zada koji metodi se pozivaju na objektu imitatoru, kojim redom i sa kojim parametrima, kao i koja vrednost se očekuje kao povratna. Na taj način se mogu imitirati ponašanja kompleksnih objekata (npr. socket) i omogućiti da programer testira da li se objekat ponaša korektno sa svim različitim stanjima. To je jednostavniji postupak nego izazivanje svih situacija na stvarnom objektu.

Rad sa objektima imitatorima obično obuhvata sledeće korake:

\begin{itemize}
\tightlist
\item
  kreiranje interfejsa za klasu koju bi trebalo testirati
\item
  kreiranje klase imitatora ručno ili pomoću nekog radnog okvira:

  \begin{itemize}
  \tightlist
  \item
    C++ - \texttt{FakeIt}, \texttt{CppUMock} (unutar \texttt{CppUnit}), \texttt{GoogleMock}
  \item
    Java - \texttt{Mockito}, \texttt{JMock}, \texttt{EasyMock}, \texttt{PowerMock}
  \item
    .NET - \texttt{Moq}
  \end{itemize}
\item
  priprema koda koji će se testirati na objektu imitatoru;
\item
  pisanje testa koji će koristiti objekat imitator umesto stvarnog objekta
\end{itemize}

Unutar testa je potrebno: -- kreirati instancu objekta klase imitatora -- podesiti ponašanje i očekivanja od objekta imitatora -- pokrenuti kod koji će koristiti objekat imitator -- po izvršavanju, porediti dobijene i očekivane vrednosti (ovaj korak obično izvršava radni okvir prilikom uništavanja objekta imitatora)

\subimport{01_cpp_manual/}{index.tex}

\subimport{02_csharp_moq/}{index.tex}

\subimport{03_csharp_dbtest/}{index.tex}
